\section{Introduction}
\label{sec:intro}

Many models built around self-consistency, soft constraints, or fixed points
require a linear solve. Deep equilibrium (DEQ) models differentiate through an
equilibrium with the implicit function theorem \citep{bai2019deep};
linear-attention and delta-rule layers mix tokens with a triangular solve
\citep{yang2024parallelizing}; implicit and steady-state graph layers solve for
a fixed point. In each case, either the forward computation or the implicit
backward contains a linear solve. Dense factorization costs $\bigO(N^3)$.
On a low-treewidth sparse graph, selected inversion instead computes only the
blocks of $\inv{A}$ on the factor's filled pattern. It costs $\bigO(n)$ at fixed
block size and elimination width, or $\bigO(n^{1.5})$ on two-dimensional grids
with nested dissection \citep{george1973nested,takahashi1973,erismantinney1975},
without forming the dense inverse. This forward algorithm is classical.

\paragraph{Research question.}
Sparse direct computation removes iteration truncation, but it does not imply
greater floating-point accuracy. On the finite precision grid in
Section~\ref{sec:numerics}, selected inversion and a well-ordered dense solve
perform comparably against a 64-bit floating-point (fp64) oracle. The question
here is instead:
\begin{quote}
\emph{When does the exactness of a structured linear-algebra primitive change
what a neural network can learn or generalize, as opposed to merely how fast
or how accurately it computes a fixed quantity?}
\end{quote}

\paragraph{Hypothesis and adjudicated outcome.}
The hypothesis was that an exact, size-agnostic global solve supplies an
inductive bias: among models with comparable training fit, it selects a
solution that generalizes better. Two preregistered probes test this on
$30$ held-out seeds with matched-fit rules
(\href{https://github.com/vanshsehrawat14/gabp-sparse-inv/blob/main/archive/docs/CONFIRMATORY_RESULTS.md}
{\nolinkurl{archive/docs/CONFIRMATORY_RESULTS.md}}). In the high-difficulty cells,
the finite alternatives match the exact arm's training fit in only
$0$--$23.3\%$ of seeds, below the preregistered $60\%$ bar. The equal-fit
comparison required by the hypothesis does not arise. The first-read evidence
therefore supports an attainability interpretation, not an inductive-bias
effect. The mechanism is more specific in each probe: fixed-budget forward
access in the routed field, and a backward gradient-estimator/optimization
limitation in the DEQ.
\begin{itemize}
  \item \emph{Size extrapolation.} A $65$-parameter model containing the
        exact solve remains accurate as the routed-field grid grows, while
        larger graph neural network (GNN) and Transformer baselines approach
        the predict-the-mean
        floor. Replacing only the solve with $K$ Jacobi sweeps produces a
        monotone reach-dependent gap, while the preregistered $K=7$ arm
        does not attain the exact arm's training fit
        (Section~\ref{sec:maze}).
  \item \emph{Stiff-regime training.} On four diagnostic cells, the
        sparse-direct implicit backward agrees with a dense fp64 oracle to
        within $1.1\times10^{-12}$, whereas finite Neumann adjoints degrade
        as $\rho(J)$ approaches one. In the preregistered training study,
        the iterative-adjoint arm does not attain comparable fit in the
        stiff cells. A later fresh-seed adversarial study shows that
        residual-optimal GMRES opens comparable fit at eight matched products,
        so the broader DEQ failure is rule-sensitive rather than a generic
        finite-access result (Section~\ref{sec:deq}).
\end{itemize}
The DEQ low-difficulty control fails its relative-gap clause. The maze
extrapolation-floor comparison passes, but a separate in-distribution
\texttt{floor\_sound} diagnostic is false.

\paragraph{Contribution.}
The exact forward engine is classical
\citep{takahashi1973,erismantinney1975,lin2011selinv,weiss2001correctness}.
Its reverse schedule also follows established selected-inverse differentiation
and cheap-gradient results \citep{zhu2019sparse,dwyer1948symbolic,baur1983complexity}.
The contributions are:
\begin{enumerate}
  \item two preregistered matched-fit studies that distinguish an
        equal-fit preference from a failure to attain comparable training
        fit;
  \item a source-generated access certificate for the routed-field probe,
        including an instance-specific label-only loss floor and an explicit
        boundary at nonlocal transforms and preconditioners;
  \item a post-hoc 20-seed residual-optimal DEQ falsifier that separates a
        finite Neumann failure from what the same operator-product budget can
        attain with GMRES; and
  \item the reference PyTorch operator family used to run the studies,
        covering differentiable selected inversion and its solve siblings.
\end{enumerate}

\paragraph{Scope.}
The evidence comes from controlled central processing unit (CPU) experiments:
routing grids up to $48\times48$, confirmatory DEQ fixed points at $n=16$, and
the exploratory GMRES adversary at $n=256$. The structural win is confined to
low-treewidth problems; frontal growth can make sparse direct impractical in
general. Benchmark-scale claims require separate protocols and matched
accuracy/cost comparisons (Section~\ref{sec:program}).
